\documentclass[10pt,a4paper]{amsart}
\usepackage[margin=19mm]{geometry}
\usepackage{amsmath,amssymb,mathtools}
\usepackage[scr=rsfs,cal=euler]{mathalfa}
\usepackage{xfrac}
\usepackage[unicode,hidelinks,bookmarks=false]{hyperref}
\newcommand{\ie}{i.e.\ }
\newcommand{\cf}{cf.\ }
\newcommand{\resp}{respectively\ }
\newcommand{\Subsection}{\subsection}
\providecommand{\nobreakdash}{\nobreak\hbox{-}}
\setlength{\emergencystretch}{2em}
\multlinegap0pt
\usepackage{bm}
\usepackage{bbm}
\usepackage{dsfont}
\usepackage{xspace}
\usepackage{cancel}
\usepackage{mathtools}
\usepackage{cleveref}
\usepackage{amsthm}
\makeatletter
\@ifundefined{thm}{\newtheorem{thm}{Thm}}{}
\@ifundefined{claim}{\newtheorem{claim}{Claim}}{}
\@ifundefined{lem}{\newtheorem{lem}[thm]{Lemma}}{}
\@ifundefined{obs}{\newtheorem{obs}[thm]{Observation}}{}
\makeatother
\title[Every graph with no $K\_7^{\vee}$-minor is $6$-colorable]{Every graph with no $K\_7^{\vee}$-minor is $6$-colorable}
\author{Sergey Norin, Agnes Totschnig}
\date{Source: arXiv:2507.03244v1 (CC BY 4.0)}
\begin{document}
\maketitle

\noindent\emph{Source and licence.} Every graph with no $K_7^{\vee}$-minor is $6$-colorable. Version arXiv:2507.03244v1 (CC BY 4.0), distributed under the Creative Commons Attribution 4.0 International licence, as recorded in the arXiv licence metadata for this exact version and shown on the abstract page https://arxiv.org/abs/2507.03244v1. Full rights notice: licenses/arxiv-2507.03244-CC-BY-4.0.txt.\par\smallskip

\noindent\textit{Editorial scope.} Counted unit: the Claim whose source label is \texttt{c:no4sep2K4} in the article (source0.tex lines 303--304, source label \texttt{c:no4sep2K4}), delivered together with the article's own complete original proof of it (source0.tex lines 306--336, one \texttt{proof} environment of 31 lines). No mathematics is written, repaired or re-derived: statement and proof are the article's own text line for line. Required context retained uncounted: o:intConnect (lines 155--160); l:K42star (lines 230--232). Every definition, display and label of those lines is retained unchanged. Omitted from this package and not redistributed: 1--154, 161--229, 233--302, 305--305, 337--694. Nothing inside a retained excerpt is omitted or altered: every retained body is delivered exactly as the article's lines are written, so no omission receipt of any kind accompanies this package. Citations: the retained excerpts carry no citation command at all, so no bibliography entry of the article is transcribed and no reference is redistributed. Reference adaptation: none was needed and none was made; every reference inside the delivered unit resolves inside it, so no unresolved reference or citation is produced. Printed slips kept literal: no printed word, symbol, hypothesis or number of the counted statement or of its proof was altered. Layout and class: the article's own class is \texttt{amsart} with packages including amssymb, amsthm, amsmath, color, fullpage, url, todonotes, comment, cleveref, tikz; the wrapper is the lane's pinned \texttt{amsart} editorial wrapper at 10pt on A4, which restates the theorem-like environments the retained text needs and adds the article's own macro definitions that the retained text needs (none), with extra wrapper packages (bm, bbm, dsfont, xspace, cancel, mathtools, cleveref). The delivered numbering is the unit's own. Assets: no retained span contains a figure environment, a graphics-inclusion command, a PDF-page inclusion, a table or a TikZ picture; no image file is copied into this package, the package declares no asset dependency and no image file is redistributed with it. Licence: version arXiv:2507.03244v1 of this article is distributed under the Creative Commons Attribution 4.0 International licence, as recorded in the arXiv licence metadata for this exact version and shown on the abstract page https://arxiv.org/abs/2507.03244v1. A full rights notice is delivered at licenses/arxiv-2507.03244-CC-BY-4.0.txt. Characters: every retained excerpt is pure ASCII, so the delivered engine, XeLaTeX with its default Latin Modern text font, reports no missing character.
\begin{obs}\label{o:intConnect} Let $G$ be a graph and let  $k \geq 0$ be an integer.	
	 \begin{enumerate}  
		\item Suppose that $|V(G)| \geq k+1$ and let $Z \subseteq V(G)$. Then $(G,Z)$ is internally  $k$-connected if and only if the graph obtained from $G$ by adding edges between all pairs of non-adjacent vertices of $Z$ is $k$-connected.
		\item If $G$ is $k$-connected and $(A,B)$ is a separation of $G$ of order at least $k$ then $(G[B],A~\cap~B)$ is internally $k$-connected. 
	\end{enumerate}	
\end{obs}

 \begin{lem}\label{l:K42star}
	Let $G$ be a graph, and let $Z \subseteq V(G)$ with $|Z|=4$ be such that $(G,Z)$ is internally $4$-connected. If $G$ has no $Z$-rooted $K^*_{4,2}$-model then $$|E(G)| \leq 4|V(G)|-10.$$ 
\end{lem}

\begin{claim}\label{c:no4sep2K4} There exists no non-trivial $4$-separation $(A,B)$ of $G$ such that $G[A]$ and $G[B]$ both have an $(A \cap B)$-rooted $K_4$-minor.
\end{claim}	

	\begin{proof} Suppose that such a separation $(A,B)$ exists. Let $k=|E(G[A \cap B])|$. Let $G_A$ and $G_B$ be the graphs obtained from $G[A]$ and $G[B]$, respectively, by adding $6-k$ missing edges of $G[A \cap B]$. Thus $A \cap B$ is a clique in $G_A$ and $G_B$.

By Observation \ref{o:intConnect} both $G_A$ and $G_B$ are $4$-connected. 	
		
As $G_A$ and $G_B$ are proper minors of $G$, neither of them is an enemy by the choice of $G$. 	Thus for $X \in \{A,B\}$ we have $|E(G_X)| \leq 4|X| - 9+\chi_X$ , where  $\chi_X =1$ if $G_X \simeq K_{2,2,2,2}$  and $\chi_X=0$, otherwise.
It follows that   \begin{align*}
	|E(G)| &= |E(G[A])| + |E(G[B])| -k = |E(G_A)|+|E(G_B)| -2(6-k)-k \\
	&\leq 4|A| - 9  + 4|B| - 9 - 12 + k + \chi_A +\chi_B 
	= 4(|V(G)| + 4) - 30 + k  + \chi_A +\chi_B \\
	&\leq 4|V(G)| - 14 + k  + \chi_A +\chi_B,
\end{align*} 
As $E(G) \geq  4|V(G)| - 8,$ it follows that $k+\chi_A +\chi_B \geq 6$ and so \begin{itemize} \item either $k = 6$, or \item $k=5$ and either  $G_A \simeq K_{2,2,2,2}$ or $G_B \simeq K_{2,2,2,2}$
 or \item $k=4$ and  $G_A \simeq K_{2,2,2,2}$ and $G_B \simeq K_{2,2,2,2}$.
 \end{itemize}
 \noindent {\bf Case 1.1: $k=6$.}
 Note that in this case neither $G[A]$ nor $G[B]$ have an $(A \cap B)$-rooted $K^*_{4,2}$-minor.
 Indeed, suppose without loss of generality that $G[A]$ has such a minor. Thus we can contract the edges of $G[A]$ to obtain a $K_6$ with four of the vertices corresponding to $A \cap B$.  Contracting any connected component of $B-A$ to a single vertex and deleting the remaining vertices of $B-A$ we obtain $K_7^{\vee}$ as a minor of $G$, a contradiction.
 
  It follows from \cref{l:K42star} that $|E(G[X])| \leq 4|X|- 10$ for $X \in \{A,B\}$ and so $$
  	|E(G)| = |E(G[A])| + |E(G[B])| -6 
  	\leq 4|A|  + 4|B| - 26 
  = 4|V(G)| - 10,
$$
a contradiction, finishing the argument in this case.

\vskip 5pt \noindent {\bf Case 1.2: $k=5$ and either  $G_A \simeq K_{2,2,2,2}$ or $G_B \simeq K_{2,2,2,2}$.}
 Assume without loss of generality that    $G_A \simeq K_{2,2,2,2}$. Let $A \cap B = \{v_1,v_2,v_3,v_4\}$ where $v_1v_2 \not \in E(G)$ and all the other pairs of vertices in $A \cap B$ are adjacent. Let $A - B = \{u_1,u_2,u_3,u_4\}$, where $v_i$ is the unique non-neighbor of $u_i$ in $G[A]$. Let $G^{+}_A$ be the graph obtained from $G$ by contracting any connected component of $B-A$ to a single vertex and deleting the remaining vertices of $B-A$, as we did in the previous case. Then $G^{+}_A$ is obtained from $G[A]$ by adding a vertex adjacent to all the vertices in  $A \cap B$. Contracting the edges $u_3v_1$ and $u_4v_2$ in $G^{+}_A$ we obtain $K_7^{\vee}$ as a minor of $G$. This contradiction finishes this case.
 
 \vskip 5pt \noindent {\bf Case 1.3: $k=4$, $G_A \simeq K_{2,2,2,2}$ and $G_B \simeq K_{2,2,2,2}$.}
  As in the previous case let $A \cap B = \{v_1,v_2,v_3,v_4\}$, and now let $B-A = \{w_1,w_2,w_3,w_4\}$ where $v_i $ is the unique non-neighbor of $w_i$ in $G[B]$. Assume without loss of generality that $v_3v_4 \in E(G)$. Contracting the edges $v_1w_2,v_2w_1$ and $w_3w_4$, we replace $G[B]$ by a clique of size five, with four vertices corresponding to $A \cap B$. It now follows as in the previous case that $G$ has a $K_7^{\vee}$ minor, again a contradiction, finishing the proof of this case and the claim.
 \end{proof}

\end{document}
